\section*{Capítulo \ref{ch:b2:unicellular-diversity} --- Diversidade dos organismos unicelulares}

\begin{solution}{exo:b2:unicellular-diversity:1}
\omterm{def:b2:unicellular-diversity:unicellular}{Unicelular}: uma célula realiza todas as funções e vive sozinha;
\omterm{def:b2:unicellular-diversity:unicellular}{colonial}: células idênticas permanecem unidas, mas cada uma poderia
viver sozinha; \omterm{def:b2:unicellular-diversity:unicellular}{multicelular}: vários tipos celulares, interdependentes,
e só a linhagem germinativa se reproduz. \emph{E.~coli} e a levedura
são \omterm{def:b2:unicellular-diversity:unicellular}{unicelulares} (o broto se destaca e vive sozinho); \emph{Gonium} é
\omterm{def:b2:unicellular-diversity:unicellular}{colonial} (16 células idênticas, cada uma capaz de fundar uma colônia);
\emph{Volvox} é \omterm{def:b2:unicellular-diversity:unicellular}{multicelular} (células somáticas mortais, poucas
células germinativas).
\end{solution}

\begin{solution}{exo:b2:unicellular-diversity:2}
$S/V = 3/r$: \qty{6}{\micro m^{-1}} para o coco,
\qty{0.06}{\micro m^{-1}} para o \omterm{def:b2:unicellular-diversity:protist}{protista}, uma diferença de cem vezes.
O coco respira mais depressa por unidade de volume: cada parte de seu
citoplasma está a uma fração de micrômetro da superfície por onde
entram o oxigênio e o alimento, ao passo que o interior do \omterm{def:b2:unicellular-diversity:protist}{protista} é
abastecido por uma superfície cem vezes menor por unidade de volume.
\end{solution}

\begin{solution}{exo:b2:unicellular-diversity:3}
Ausência de núcleo (um \omterm{def:b2:unicellular-diversity:prokaryote}{nucleoide}) contra um envoltório nuclear;
ausência de mitocôndrias e de endomembranas, com a cadeia respiratória
na membrana plasmática, contra organelas; um flagelo que é uma hélice
rígida de flagelina girada por prótons, contra um cílio $9+2$ que se
curva pela ação da dineína e do ATP; e ainda o tamanho
(\qty{1}{\micro m} contra dezenas) e uma parede de peptidoglicano. Uma
bactéria não consegue fagocitar: sua parede a impede de englobar uma
partícula, o que uma ameba faz com seus pseudópodes.
\end{solution}

\begin{solution}{exo:b2:unicellular-diversity:4}
Cílios: locomoção e condução do alimento para a boca; sulco oral:
canaliza as partículas para a citofaringe; vacúolo digestivo: digere
as bactérias englobadas com enzimas lisossômicas; \omterm{prop:b2:unicellular-diversity:osmoregulation}{vacúolo contrátil}:
expulsa a água que entra por osmose; macronúcleo: o núcleo funcional,
com centenas de cópias de genes transcritas para a vida diária;
micronúcleo: o núcleo diploide germinativo, usado na meiose e na
conjugação.
\end{solution}

\begin{solution}{exo:b2:unicellular-diversity:5}
$t = x^2/2D$: $(10^{-6})^2/10^{-9} = \qty{1}{ms}$ para a bactéria;
$(5\times 10^{-4})^2/10^{-9} = \qty{250}{s}$, cerca de quatro minutos,
para a ameba. A ameba não pode contar com a difusão para distribuir
os metabólitos: precisa de correntes citoplasmáticas e de transporte
por motores ao longo do citoesqueleto, dos quais a bactéria prescinde.
\end{solution}

\begin{solution}{exo:b2:unicellular-diversity:6}
$Re = \rho vL/\eta$: \emph{Paramecium} $10^3\times 10^{-3}\times
2\times 10^{-4}/10^{-3} = 0.2$; bactéria $10^3\times 2\times
10^{-5}\times 2\times 10^{-6}/10^{-3} = 4\times 10^{-5}$; nadador
$10^3\times 1\times 2/10^{-3} = 2\times 10^6$. Com $Re \gg 1$, a inércia
leva o nadador adiante após a braçada; com $Re \ll 1$, o arrasto
viscoso detém a bactéria em menos de um nanômetro, de modo que ela só
se move enquanto bate os flagelos.
\end{solution}

\begin{solution}{exo:b2:unicellular-diversity:7}
$v_s = 2r^2\Delta\rho g/9\eta = 2\times 10^{-10}\times 100\times
9.81/(9\times 10^{-3}) = \qty{2.2e-5}{m/s}$; \qty{50}{m} levam
$50/2.2\times 10^{-5} = \qty{2.3e6}{s}$, 27 dias. Reduzir o raio à
metade divide $v_s$ por 4: 108 dias. Reduzir $\Delta\rho$ a
\qty{20}{kg/m^3} a divide por 5: 135 dias.
\end{solution}

\begin{solution}{exo:b2:unicellular-diversity:8}
$V = \tfrac{4}{3}\pi\times 100\times 25\times 25 =
\qty{2.6e5}{\micro m^3}$. Em \qty{1800}{s}: $2.6\times 10^5/1800 =
\qty{145}{\micro m^3/s}$; como $\qty{1}{\micro m^3} =
\qty{1e-15}{L}$, isso dá \qty{1.5e-13}{L/s}.
\end{solution}

\begin{solution}{exo:b2:unicellular-diversity:9}
Uma duplicação ao longo de \qty{1}{mm} é um aumento de cerca de
\qty{0.07}{\%} por
micrômetro (pois $2^{1/1000} - 1 = 0.0007$), logo \qty{0.07}{\%}
de um lado a outro da célula. Uma corrida de \qty{1}{s} percorre
\qty{20}{\micro m}: \qty{1.4}{\%}. Para resolver uma diferença de
\qty{0.07}{\%}, a célula teria de contar cerca de
$(1/0.0007)^2 = 2\times 10^6$ moléculas em cada
extremidade, muito mais do que as que se ligam a seus poucos milhares
de receptores em um segundo; \qty{1.4}{\%} exige cerca de \num{5000},
o que é viável. A comparação temporal transforma uma medida espacial
inviável numa medida possível, deixando que a corrida faça a
integração.
\end{solution}

\begin{solution}{exo:b2:unicellular-diversity:10}
Superfície $4\pi r^2$ com $r = \qty{250}{\micro m}$: \qty{7.9e5}{\micro
m^2}; volume citoplasmático $\approx$ superfície $\times$ \qty{2}{\micro m}
$= \qty{1.6e6}{\micro m^3}$; razão \qty{0.5}{\micro m^{-1}}. Para
\emph{E.~coli}, um cilindro: $S/V = 2\pi r L/\pi r^2 L = 2/r =
\qty{4}{\micro m^{-1}}$. A gigante é apenas oito vezes pior que
\emph{E.~coli}, e não 500 vezes, porque cada ponto de seu citoplasma
está a menos de \qty{2}{\micro m} da superfície. O vacúolo armazena
nitrato, seu aceptor de elétrons, para que ela possa respirar sulfeto
no sedimento anóxico entre as raras agitações que trazem nitrato
(\cref{ch:b2:microbial-metabolism}).
\end{solution}

\begin{solution}{exo:b2:unicellular-diversity:11}
Uma célula não consegue nadar e se dividir ao mesmo tempo. Uma colônia
de 16 células que para de nadar durante as horas da divisão afunda,
pela \omterm{thm:b2:unicellular-diversity:stokes}{lei de Stokes}, uma distância desprezível (seu raio é de algumas
dezenas de micrômetros, sua velocidade de sedimentação de micrômetros
por segundo); a colônia inteira pode se dividir sem perder nada. Uma
esfera de \num{10000} células com raio de
\qty{500}{\micro m} afundaria cem vezes mais depressa ---
milímetros por segundo, metros em uma hora --- para fora da luz:
compensa manter a maioria das células nadando enquanto algumas se
dividem. As células germinativas são grandes porque precisam levar os
recursos para construir, por divisões sucessivas e sem se alimentar,
uma esfera-filha inteira de milhares de células; por isso são
abastecidas pelo soma.
\end{solution}

\begin{solution}{exo:b2:unicellular-diversity:12}
Um clado é um ancestral e todos os seus descendentes; um grau é um
nível de organização alcançado por várias linhagens. Os \omterm{def:b2:unicellular-diversity:protist}{protistas} ---
algas verdes (com as plantas terrestres), diatomáceas (com as algas
pardas), ciliados (com os dinoflagelados e o \emph{Plasmodium}),
amebas (perto dos animais e dos fungos) --- têm em comum apenas a
célula eucariótica vivendo sozinha, e seu último ancestral comum é o
de todos os eucariotos, nós incluídos: um grau. ``\omterm{def:b2:unicellular-diversity:prokaryote}{Procarioto}'' reúne
bactérias e arqueias por aquilo que lhes falta; as arqueias são mais
próximas dos eucariotos do que das bactérias, de modo que o termo é,
de novo, um grau. As cianobactérias descendem de um único ancestral
que inventou a fotossíntese oxigênica e incluem todos os seus
descendentes (deixando de lado a linhagem do cloroplasto): um clado.
\end{solution}

\begin{solution}{pb:b2:unicellular-diversity:1}
\textbf{1.} $1\times 1\times 0.1 = \qty{0.1}{mm^3} = \qty{0.1}{\micro L}$.
\textbf{2.} $24/\qty{0.1}{\micro L} = 240$ por \unit{\micro L} $=
\qty{2.4e5}{}$ células por mililitro.
\textbf{3.} $V = \tfrac{4}{3}\pi\times 125 = \qty{524}{\micro m^3}$;
$S/V = 3/5 = \qty{0.6}{\micro m^{-1}}$.
\textbf{4.} $2.4\times 10^5\times 524 = \qty{1.26e8}{\micro m^3}$ por
mililitro, e $\qty{1}{mL} = \qty{1e12}{\micro m^3}$: uma fração
$1.3\times 10^{-4}$, \qty{0.013}{\%}.
\textbf{5.} $\qty{1000}{m^3} = \qty{1e9}{mL}$: $2.4\times 10^{14}$
células.
\textbf{6.} A cianobactéria é um filamento de células muito menores
(alguns micrômetros), verde-azuladas, sem cloroplasto nem núcleo
distintos, e não nada com flagelos; a alga é uma célula ovoide isolada
com dois flagelos, um estigma e um único cloroplasto em forma de taça.
\textbf{7.} $384/24 = 16 = 2^4$: quatro duplicações em \qty{48}{h},
$T = \qty{12}{h}$.
\textbf{8.} $k = \ln 2/T = 0.693/12 = \qty{0.058}{h^{-1}}$.
\textbf{9.} $t = \ln(10^8/2.4\times 10^5)/k = \ln(417)/0.058 =
6.03/0.058 = \qty{104}{h}$, cerca de 4.3 dias.
\textbf{10.} Os nutrientes (fosfato, nitrato) se esgotam; as células
se sombreiam umas às outras, de modo que a luz por célula diminui; os
herbívoros (ciliados, rotíferos, \emph{Daphnia}) se multiplicam à
custa delas; a \qty{1e8}{} por mililitro, as células ocupariam cerca de
cinco por cento do volume e esgotariam o $\mathrm{CO_2}$
dissolvido.
\textbf{11.} Crescimento em apenas metade do tempo: a taxa média cai à
metade, \qty{208}{h}, 8.7 dias.
\textbf{12.} A contagem mede a população, não células individuais:
se uma célula libera 8 filhas a cada \qty{36}{h}, a população é
multiplicada por 8 $= 2^3$ em três \omterm{def:b2:unicellular-diversity:division}{tempos de geração} de \qty{12}{h}
cada. O \omterm{def:b2:unicellular-diversity:division}{tempo de geração} é o tempo médio de duplicação do número,
qualquer que seja o modo como as divisões se agrupam.
\textbf{13.} $Re = 10^3\times 10^{-4}\times 10^{-5}/10^{-3} = 10^{-3}$.
\textbf{14.} $m = 1050\times 5.24\times 10^{-16} = \qty{5.5e-13}{kg}$;
$6\pi\eta r = 6\pi\times 10^{-3}\times 5\times 10^{-6} =
\qty{9.4e-8}{kg/s}$; $d = 5.5\times 10^{-13}\times 10^{-4}/9.4\times
10^{-8} = \qty{5.8e-10}{m}$: meio nanômetro.
\textbf{15.} $v_s = 2\times 25\times 10^{-12}\times 50\times
9.81/(9\times 10^{-3}) = \qty{2.7e-6}{m/s}$, \qty{2.7}{\micro m/s}.
\textbf{16.} Afundar \qty{1}{m}: $1/2.7\times 10^{-6} =
\qty{3.7e5}{s}$, 4.3 dias; subir nadando a \qty{100}{\micro m/s}:
\qty{1e4}{s}, 2.8 horas. O nado vence por um fator de 37.
\textbf{17.} $c_\infty = \qty{1e-2}{mol/m^3}$: $J = 4\pi\times
1.9\times 10^{-9}\times 5\times 10^{-6}\times 10^{-2} =
\qty{1.2e-15}{mol/s} = 7.2\times 10^8$ moléculas por segundo.
\textbf{18.} \qty{100}{pg} $= \qty{1e-10}{g} = \qty{8.3e-12}{mol}$ de
carbono, duplicados em \qty{12}{h} $= \qty{43200}{s}$:
\qty{1.9e-16}{mol/s} $= 1.2\times 10^8$ átomos por segundo. Suprimento
sobre demanda: $7.2/1.2 = 6$.
\textbf{19.} O suprimento cresce dez vezes ($\propto r$), a demanda
mil vezes ($\propto r^3$): a razão passa a $6/100 = 0.06$; tal célula
só poderia crescer a seis por cento da taxa, de modo que uma alga
desse tamanho alimentada por difusão é impossível sem agitação,
transporte interno ou um metabolismo muito mais lento.
\textbf{20.} $\Delta\Pi = RT\Delta c = 8.314\times 293\times 99 =
\qty{2.4e5}{Pa}$, \qty{2.4}{bar}.
\textbf{21.} Área $4\pi(5\times 10^{-6})^2 = \qty{3.14e-10}{m^2}$;
influxo $L_p A\Delta\Pi = 10^{-14}\times 3.14\times 10^{-10}\times
2.4\times 10^5 = \qty{7.6e-19}{m^3/s} = \qty{0.76}{\micro m^3/s}$.
\textbf{22.} $524/0.76 = \qty{690}{s}$, cerca de onze minutos.
\textbf{23.} $P = \Delta\Pi\times$ vazão $= 2.4\times 10^5\times
7.6\times 10^{-19} = \qty{1.8e-13}{W}$.
\textbf{24.} Um ATP: $5\times 10^4/6.02\times 10^{23} =
\qty{8.3e-20}{J}$: $1.8\times 10^{-13}/8.3\times 10^{-20} = 2.2\times
10^6$ ATP por segundo. Fixação de carbono: $3\times 1.2\times 10^8 =
3.6\times 10^8$ ATP por segundo. O vacúolo custa \qty{0.6}{\%} disso
--- um seguro barato contra o rompimento.
\textbf{25.} Influxo de água de \qty{0.76}{\micro m^3/s} (o volume da
célula em onze minutos); sem \omterm{prop:b2:unicellular-diversity:osmoregulation}{vacúolo contrátil}, a célula dobraria de
volume em cerca de \qty{700}{s}; manter-se sem inchar custa pelo menos
$2\times 10^6$ ATP por segundo, menos de um por cento do que sua
fotossíntese gasta com o carbono.
\end{solution}
